\section{TRABALHOS RELACIONADOS}

\subsection{Modelagem Temporal com Deep Learning em Hidrologia}

A modelagem de qualidade da água tem avançado de abordagens clássicas para redes neurais profundas,
com destaque para modelos recorrentes em séries temporais ambientais. \citeonline{pyo2023} mostram
ganhos consistentes com LSTM em ambientes aquáticos interiores, enquanto \citeonline{qi2020} e
\citeonline{liu2019} reforçam a efetividade de arquiteturas recorrentes e IoT para previsão de
parâmetros fisico-químicos. Em nível de síntese, revisões recentes confirmam a predominância de
Deep Learning para padrões não lineares em hidrologia \cite{yan2024,pan2025}. Apesar da precisão,
esses trabalhos também indicam custo computacional e energético elevados para execução contínua em borda.

\subsection{Arquiteturas Híbridas e Pré-processamento de Sinal}

Abordagens híbridas são usadas para melhorar robustez em cenários ruidosos. \citeonline{wu2022}
integram ANN, Wavelet e LSTM para elevar desempenho preditivo, enquanto \citeonline{ahmed2019},
\citeonline{kalaivanan2022} e \citeonline{aladhaileh2021} comparam múltiplos algoritmos de
aprendizado de máquina para parâmetros de qualidade da água. Esses resultados mostram ganhos de
acurácia, porém com maior custo de pré-processamento e inferência quando o objetivo é execução
local em hardware restrito.

\subsection{Avanços em Redes Neurais Bioinspiradas}

Entre os modelos bioinspirados recentes, \citeonline{chadalavada2025} propõem Gated Liquid Neural
Networks para índice de qualidade da água, com resultados promissores de predição e classificação.
Embora eficientes, tais modelos ainda operam em regime contínuo. Em contraste, SNNs exploram
processamento orientado a eventos (\textit{event-driven}), potencialmente mais compatível com
restrições energéticas de sensores de borda.

\subsection{Detecção de Anomalias e Confiabilidade de Sensores}

A detecção de anomalias e a confiabilidade da instrumentação são pilares em monitoramento contínuo.
\citeonline{elshafeiy2023} apresentam abordagem multivariada em tempo real para detecção de eventos
anômalos, enquanto \citeonline{lamrini2024} avaliam modelos clássicos para detecção de outliers em
sensoriamento inteligente. Complementarmente, revisões de sistemas IoT para água mostram que o
desempenho analítico depende de arquitetura de dados, calibração e robustez da rede de sensores
\cite{zulkifli2022,miller2025}.

\subsection{Edge Computing e Estratégias de Eficiência Energética}

Edge AI tem sido proposta para reduzir latência e dependência de nuvem em monitoramento hídrico
\cite{rahu2023,lowe2022}. Estudos aplicados em redes de sensores sem fio e monitoramento remoto
mostram viabilidade operacional, mas também destacam limites de hardware, energia e conectividade
\cite{rosero2020,alprol2024}. No contexto amazônico, essas limitações são críticas e motivam
arquiteturas de inferência local com menor custo computacional.
Além disso, a literatura de visão estratégica sobre monitoramento rural enfatiza que arquiteturas
neuromórficas na borda podem reduzir consumo energético e tornar os sistemas mais resilientes em
áreas com infraestrutura limitada \cite{aral2024}.
Resultados recentes em IIoT embarcado também reportam ganhos energéticos com hardware
neuromórfico em comparação com plataformas convencionais \cite{azadi2025}, enquanto soluções
baseadas em ESP32 mostram viabilidade para monitoramento em tempo real em infraestrutura crítica
\cite{alshareeda2025}.

\subsection{Fundamentos de Redes Neurais Pulsadas (SNN) em Computação Neuromórfica}

Para mitigar o custo energético de ANNs convencionais em inferência contínua, a computação
neuromórfica propõe processamento por eventos. Em SNNs, a informação é codificada em pulsos
assíncronos (\textit{spikes}), reduzindo atividade computacional quando não há variação relevante
nos sinais.
No nível de hardware, aceleradores SNN com \textit{compute-in-memory} já reportam ganhos de
desempenho e eficiência energética em tecnologias avançadas de fabricação, reforçando a viabilidade
de arquiteturas neuromórficas para cenários de borda \cite{huijbregts2024}. Em paralelo, revisões
recentes sobre SNNs consolidam fundamentos de modelagem, treinamento e métricas de latência e
energia, apoiando sua adoção em aplicações sensíveis a consumo e tempo de resposta
\cite{aribe2025future}.

Entre os diversos modelos que descrevem a dinâmica do neurônio pulsado, o Leaky Integrate-and-Fire 
(LIF) destaca-se pelo equilíbrio entre plausibilidade biológica e eficiência computacional. A dinâmica 
do potencial de membrana $V(t)$ neste modelo é governada pela seguinte equação diferencial linear:

$$\tau_m \frac{dV(t)}{dt} = -(V(t) - V_{rest}) + R_m I(t)$$

Nesta formulação, $\tau_m$ representa a constante de tempo da membrana, que define a taxa de decaimento 
do potencial na ausência de estímulo; $V_{rest}$ é o potencial de repouso para o qual o neurônio tende 
a retornar; $R_m$ é a resistência da membrana; e $I(t)$ é a corrente de entrada sináptica. O mecanismo 
de disparo ocorre quando $V(t)$ ultrapassa um limiar predefinido $V_{th}$, momento em que um spike é 
emitido e o potencial é instantaneamente redefinido para um valor de reset ($V_{reset}$).

Para implementação em sistemas digitais discretos, como microcontroladores, a equação contínua é 
discretizada, resultando na seguinte aproximação para o passo de tempo $t+1$:

$$V[t+1] = V[t] + \frac{1}{\tau_m} (-(V[t] - V_{rest}) + I[t]) \Delta t$$

Esta formulação evidencia a vantagem potencial das SNNs em borda: na ausência de entrada significativa
($I[t] \approx 0$), o sistema evolui com baixa atividade, favorecendo eficiência energética em
dispositivos com recursos limitados.

\subsection{Contexto local e lacuna de pesquisa}

No contexto do Lago do Serpa, estudos locais sobre território e ambiente reforçam a relevância social
e ecológica de monitoramento contínuo na região \cite{nogueira2022}. Em paralelo, a literatura clássica
de previsão de qualidade da água com ANN já demonstrava potencial de modelagem em diferentes cenários
\cite{palani2008}, mas sem foco em inferência neuromórfica embarcada. Assim, a lacuna abordada neste
trabalho está na integração de predição multimodal, mapeamento espacial e implantação offline em
dispositivo móvel, mantendo viabilidade para uso em borda.

\subsection{Tabela comparativa dos estudos e posicionamento do trabalho}

A Tabela~\ref{tab:comparativo_estado_arte} sintetiza estudos representativos e posiciona
o presente trabalho em termos de tipo de modelo, escopo de saída e foco em implantação
de borda.

\begin{table*}[t]
\centering
\caption{Comparativo entre estudos da literatura e o presente trabalho}
\label{tab:comparativo_estado_arte}
\footnotesize
\setlength{\tabcolsep}{4pt}
\renewcommand{\arraystretch}{1.15}
\resizebox{\textwidth}{!}{%
\begin{tabular}{p{2.8cm} p{2.4cm} p{2.7cm} p{2.2cm} p{2.1cm} p{2.1cm} p{2.1cm}}
\toprule
\textbf{Estudo} & \textbf{Abordagem} & \textbf{Dados/escopo} & \textbf{Saída} & \textbf{Edge/offline} & \textbf{Mapa espacial} & \textbf{Contribuição central} \\
\midrule
\citeonline{palani2008} & ANN clássica & Águas costeiras de Singapura; série temporal & Variáveis físico-químicas (salinidade, temperatura, OD, clorofila-a) & Não reportado & Não & Previsão de variáveis de qualidade com ANN \\
\citeonline{qi2020} & LSTM-RNN + sensoriamento remoto & Lago Taihu; dados in situ + imagens de satélite & 4 parâmetros (pH, OD, CODMn, NH\textsubscript{3}-N) & Não reportado & Sim & Inversão de parâmetros e geração de mapas temáticos \\
\citeonline{pyo2023} & Revisão de LSTM (survey) & Revisão de estudos em águas interiores & Não aplicável (artigo de revisão) & Não aplicável & Não aplicável & Síntese crítica de arquiteturas LSTM e extensões \\
\citeonline{rahu2023} & IoT + ML (framework aplicado) & Canal Rohri (Paquistão), sensores IoT + modelos ML & WQI (regressão) e WQC (classificação) & Parcial (IoT/edge discutido) & Não & Framework integrado de aquisição e predição \\
\citeonline{rosero2020} & WSN + ML supervisionado & Rio Tahuando (Equador), rede de sensores & Classes de contaminação & Não (arquitetura conectada) & Sim & Sistema WSN com interface e mapa de contaminação \\
\citeonline{chadalavada2025} & Gated Liquid Neural Networks & Base pública da Índia (2005--2014) & WQI (regressão) + classes de qualidade & Não reportado & Não & Modelo bioinspirado para predição e classificação \\
\citeonline{aral2024} & Position paper (Neuromorphic Edge AI) & Monitoramento ambiental rural/remoto & Conceitual (agenda de pesquisa) & Sim (proposição) & Não aplicável & Direções para Edge AI neuromórfica em áreas rurais (sem implementação experimental) \\
\textbf{Este trabalho} & \textbf{SNN multimodal + app móvel} & \textbf{Lago do Serpa (AM), entradas: coordenadas e variáveis físico-químicas} & \textbf{Predição de pH, OD e turbidez + índice de poluição composto} & \textbf{Sim (inferência offline no celular)} & \textbf{Sim (mapa e heatmap)} & \textbf{Integra modelagem neuromórfica, análise espacial e implantação prática} \\
\bottomrule
\end{tabular}
}
\end{table*}

%----------------------------------------------------------------------------------------
%   MÉTODOS
%----------------------------------------------------------------------------------------
